\section{Constructive Context Uncertainty}
\label{app:refinement}

Fix a query $Q=[a,b]$, a complete strict ranking, and a budget $k$.
Let $P$ and $H$ be the exact possible and guaranteed support sets.
Let $F$ contain any source items retained under a fixed fallback policy.
Fallback items have unknown temporal support.
They enter every operational retrieval mask.
Define
\[
 U=\operatorname{Top}_k(P\cup F)\setminus(H\cup F).
\]
The exact context theorem makes $U$ empty precisely when the context is stable.

\paragraph{Construct two witnessing timelines.}
Choose $c\in U$.
Fewer than $k$ items in $P\cup F$ outrank $c$.
Every world supporting $c$ must therefore select it.
Set $t=\max(a,\ell_c)$ and $x_c=\min(u_c,t)$.
Possible support gives $t\le b$ and $t<p_c$.
For every nonpivot event $d$, assign
\[
 x_d=\begin{cases}
 \ell_d,&\ell_d\le x_c,\\
 u_d,&\ell_d>x_c.
 \end{cases}
\]
The first case places $d$ no later than $c$.
For every accepted witness, the second case places $d$ after $t$.
Indeed, if $x_c=t$, then $u_d\ge \ell_d>t$.
Otherwise, $x_c=u_c<t$ and $\ell_d>u_c$.
The definition of $p_c$ then gives $u_d\ge p_c>t$.
This rule also makes valid choices for all events outside the incoming witness set.
Those choices cannot affect whether $c$ is active.
The assignments respect all independent date intervals.
Claim $c$ is active at $t$, so this world selects it.

For a second world, first consider $u_c>b$.
Set $x_c=u_c$ and every other date to its lower bound.
The claim starts after the query and cannot enter its context.
Otherwise, failed guaranteed support gives $a>\ell_c$ and $r_c\le a$.
Choose the stored witness $d$ for $r_c$.
Set $x_c=\ell_c$ and $x_d=\min(u_d,a)$.
The witness conditions give $\ell_d\le x_d\le u_d$ and $\ell_c<x_d\le a$.
Thus $d$ retires $c$ by the query start.
Set every remaining date to its lower bound.
This world cannot select $c$.
The two contexts differ.

Let $n$ count modeled events and $N$ count ranked excerpts, including fixed fallback items.
The construction scans the event bounds and pivot certificate in $O(n)$ time.
It returns $O(n)$ date values and needs no retained pair list.
The pivot's rank condition already proves that these assignments yield different contexts.
Optional replay scans the accepted pairs and ranking in $O(n+m+N)$ time.
Explicit end events enter the same construction but remain hidden from the context.
Direct context replay retains the accepted pairs and uses $O(n+m+N)$ audit state.
The compact $O(n)$ certificate output stores support boundaries, rather than the full graph.

\paragraph{Every current pivot must be resolved.}
Narrow the date intervals without emptying any interval.
Keep all other inputs fixed.
Let $P'$ and $H'$ denote the refined support sets.
Then $P'\subseteq P$ and $H\subseteq H'$.
Take $c\in U$ that remains in $P'$.
It still has fewer than $k$ possible items above it.
Thus $c\in\operatorname{Top}_k(P'\cup F)$.
If the refined context is stable, the exact context theorem requires $c\in H'\cup F$.
Since $c\notin F$, it must belong to $H'$.
Therefore, each current pivot becomes guaranteed or impossible.

This condition is necessary but not sufficient.
For example, rank two claims $c,d$ with independent start bounds $[0,2]$.
Use no replacement pairs, point query $1$, and budget $1$.
Initially, the frontier is $\{c\}$.
Refining $c$ to start $2$ resolves that pivot but exposes frontier $\{d\}$.
The result therefore gives no minimum number of date refinements.

